\documentclass[11pt]{article}

% ============================================================
% Packages
% ============================================================
\usepackage[utf8]{inputenc}
\usepackage[margin=1in]{geometry}
\usepackage{amsmath,amssymb,amsthm}
\usepackage{graphicx}
\usepackage{booktabs}
\usepackage{xcolor}
\usepackage{hyperref}
\usepackage{tikz}
\usepackage{enumitem}
\usepackage{pgfplots}

\pgfplotsset{compat=1.18}

% ============================================================
% Document Formatting
% ============================================================
\setlength{\parskip}{0.6em}
\setlength{\parindent}{0pt}

% ============================================================
% List Formatting
% ============================================================
\setlist[enumerate]{
    leftmargin=2em,
    itemsep=0.25em
}

% ============================================================
% Solution Environment
% ============================================================
\newenvironment{solution}
    {\par\noindent\textbf{Solution.}\par}
    {\par}

% ============================================================
% Link Styling
% ============================================================
\hypersetup{
    colorlinks=true,
    linkcolor=blue,
    urlcolor=cyan
}

% ============================================================
% Plain note callouts
% ============================================================
\newenvironment{definitionbox}[1]
    {\par\medskip\noindent\textbf{Definition: #1.}\par}
    {\par\medskip}
\newenvironment{theorembox}[1]
    {\par\medskip\noindent\textbf{Theorem: #1.}\par}
    {\par\medskip}
\newenvironment{examplebox}[1]
    {\par\medskip\noindent\textbf{Example: #1.}\par}
    {\par\medskip}
\newenvironment{keyideabox}
    {\par\medskip\noindent\textbf{Key idea.}\par}
    {\par\medskip}
\newenvironment{formulabox}
    {\par\medskip\noindent\textbf{Important formula.}\par}
    {\par\medskip}

% ============================================================
% Metadata
% ============================================================
\title{Ordinary Differential Equations}
\author{Chris Ritter}
\date{September 15, 2026}

% ============================================================
% Document
% ============================================================

\begin{document}

\maketitle
\tableofcontents
\newpage

These notes introduce several important classes of first-order
differential equations, with particular emphasis on exact equations
and integrating factors.

% ============================================================
\section{First-Order Equations}
% ============================================================

Some important types of first-order differential equations are:

\begin{itemize}
    \item Linear equations
    \item Separable equations
    \item Exact equations
\end{itemize}

% ============================================================
\subsection{Exact Differential Equations}
% ============================================================

A first-order differential equation can be written in the form

\[
M(x,y)+N(x,y)\frac{dy}{dx}=0.
\]

\begin{definitionbox}{Exact Differential Equation}

The differential equation

\[
M(x,y)+N(x,y)\frac{dy}{dx}=0
\]

is called \textbf{exact} if there exists a function
$\psi=\psi(x,y)$ such that

\[
\psi_x=M,
\qquad
\psi_y=N.
\]

The function $\psi(x,y)$ is called a
\textbf{potential function}.

\end{definitionbox}

If such a function $\psi$ exists, then by the chain rule,

\[
\frac{d}{dx}\psi(x,y(x))
=
\psi_x+\psi_y\frac{dy}{dx}.
\]

Since

\[
\psi_x=M
\]

and

\[
\psi_y=N,
\]

we obtain

\[
\frac{d}{dx}\psi(x,y)
=
M(x,y)+N(x,y)\frac{dy}{dx}.
\]

Therefore, if

\[
M(x,y)+N(x,y)\frac{dy}{dx}=0,
\]

then

\[
\frac{d}{dx}\psi(x,y)=0.
\]

It follows that

\[
\boxed{\psi(x,y)=C}.
\]

\begin{keyideabox}

For an exact differential equation, the main goal is to find the
potential function $\psi(x,y)$. Once $\psi$ is found, the solution is

\[
\boxed{\psi(x,y)=C}.
\]

\end{keyideabox}

% ============================================================
\subsection{Example}
% ============================================================

\begin{examplebox}{Finding the Potential Function}

Consider

\[
2x+y^2+2xy\frac{dy}{dx}=0.
\]

Identify

\[
M(x,y)=2x+y^2,
\qquad
N(x,y)=2xy.
\]

We seek $\psi(x,y)$ satisfying

\[
\psi_x=2x+y^2,
\qquad
\psi_y=2xy.
\]

Integrating $\psi_x$ with respect to $x$ gives

\[
\psi(x,y)
=
\int(2x+y^2)\,dx+h(y),
\]

so

\[
\psi(x,y)
=
x^2+xy^2+h(y).
\]

Differentiate with respect to $y$:

\[
\psi_y
=
2xy+h'(y).
\]

Since $\psi_y=N=2xy$,

\[
2xy+h'(y)=2xy,
\]

and therefore

\[
h'(y)=0.
\]

Thus,

\[
h(y)=C.
\]

Absorbing the constant into the right-hand side gives

\[
\boxed{x^2+xy^2=C}.
\]

\end{examplebox}

% ============================================================
\subsection{Checking Whether an Equation Is Exact}
% ============================================================

The following theorem gives a convenient test for exactness.

\begin{theorembox}{Exactness Criterion}

Let $M$, $N$, $M_y$, and $N_x$ be continuous in a rectangular region
$R$. Then

\[
M(x,y)+N(x,y)\frac{dy}{dx}=0
\]

is exact if and only if

\[
\boxed{M_y=N_x}
\]

for all $(x,y)\in R$.

\end{theorembox}

To see why this condition appears, suppose that the equation is exact.
Then there exists $\psi(x,y)$ such that

\[
\psi_x=M,
\qquad
\psi_y=N.
\]

Taking mixed partial derivatives,

\[
\psi_{xy}=M_y
\]

and

\[
\psi_{yx}=N_x.
\]

Because the relevant partial derivatives are continuous,

\[
\psi_{xy}=\psi_{yx}.
\]

Hence,

\[
\boxed{M_y=N_x}.
\]

% ============================================================
\subsection{Finding the Potential Function}
% ============================================================

Suppose

\[
M(x,y)+N(x,y)\frac{dy}{dx}=0
\]

is exact.

Since

\[
\psi_x=M,
\]

we integrate $M$ with respect to $x$:

\[
\boxed{
\psi(x,y)
=
\int M(x,y)\,dx+h(y)
}.
\]

The function $h(y)$ appears because we are integrating with respect
to $x$. Any term depending only on $y$ has derivative zero with respect
to $x$.

We then differentiate with respect to $y$:

\[
\psi_y
=
\frac{\partial}{\partial y}
\left(
\int M(x,y)\,dx
\right)
+h'(y).
\]

Since $\psi_y=N$,

\[
h'(y)
=
N(x,y)
-
\frac{\partial}{\partial y}
\left(
\int M(x,y)\,dx
\right).
\]

Finally, integrate $h'(y)$ to obtain $h(y)$ and therefore $\psi(x,y)$.

% ============================================================
\subsection{Example: Checking Exactness and Solving}
% ============================================================

\begin{examplebox}{Exact Equation}

Consider

\[
\left(y\cos x+2xe^y\right)
+
\left(\sin x+x^2e^y-1\right)\frac{dy}{dx}
=0.
\]

Identify

\[
M(x,y)=y\cos x+2xe^y
\]

and

\[
N(x,y)=\sin x+x^2e^y-1.
\]

First,

\[
M_y
=
\cos x+2xe^y.
\]

Also,

\[
N_x
=
\cos x+2xe^y.
\]

Therefore,

\[
M_y=N_x,
\]

so the equation is exact.

Since

\[
\psi_x=y\cos x+2xe^y,
\]

integrate with respect to $x$:

\[
\psi(x,y)
=
y\sin x+x^2e^y+h(y).
\]

Differentiate with respect to $y$:

\[
\psi_y
=
\sin x+x^2e^y+h'(y).
\]

Set this equal to $N$:

\[
\sin x+x^2e^y+h'(y)
=
\sin x+x^2e^y-1.
\]

Thus,

\[
h'(y)=-1,
\qquad
h(y)=-y+C.
\]

Hence,

\[
\boxed{
y\sin x+x^2e^y-y=C
}.
\]

\end{examplebox}

% ============================================================
\section{Integrating Factors}
% ============================================================

An equation of the form

\[
M(x,y)+N(x,y)\frac{dy}{dx}=0
\]

may not be exact. We can sometimes multiply the equation by a
function $\mu(x,y)$ so that the resulting equation becomes exact.

\begin{definitionbox}{Integrating Factor}

An \textbf{integrating factor} is a function $\mu(x,y)$ that can be
multiplied into a differential equation so that the resulting equation
is exact.

\end{definitionbox}

We want

\[
\mu M+\mu N\frac{dy}{dx}=0
\]

to be exact. Therefore,

\[
\frac{\partial}{\partial y}(\mu M)
=
\frac{\partial}{\partial x}(\mu N).
\]

Using the product rule,

\[
\mu_yM+\mu M_y
=
\mu_xN+\mu N_x.
\]

% ============================================================
\subsection{Integrating Factor Depending Only on
\texorpdfstring{$x$}{x}}
% ============================================================

Suppose

\[
\mu=\mu(x).
\]

Then

\[
\mu_y=0.
\]

Therefore,

\[
\mu M_y
=
\mu_xN+\mu N_x.
\]

Rearranging,

\[
\mu(M_y-N_x)=\mu_xN.
\]

Thus,

\[
\frac{\mu_x}{\mu}
=
\frac{M_y-N_x}{N}.
\]

Therefore,

\[
\boxed{
\frac{d\mu}{\mu}
=
\frac{M_y-N_x}{N}\,dx
}.
\]

If the right-hand side is a function of $x$ alone, then

\[
\boxed{
\mu(x)
=
e^{\int\frac{M_y-N_x}{N}\,dx}
}.
\]

% ============================================================
\subsection{Integrating Factor Depending Only on
\texorpdfstring{$y$}{y}}
% ============================================================

Similarly, suppose

\[
\mu=\mu(y).
\]

Then

\[
\mu_x=0.
\]

Therefore,

\[
\mu_yM+\mu M_y
=
\mu N_x.
\]

Rearranging,

\[
\mu_yM
=
\mu(N_x-M_y).
\]

Thus,

\[
\frac{\mu_y}{\mu}
=
-\frac{M_y-N_x}{M}.
\]

Therefore,

\[
\boxed{
\frac{d\mu}{\mu}
=
-\frac{M_y-N_x}{M}\,dy
}.
\]

If the right-hand side is a function of $y$ alone, then

\[
\boxed{
\mu(y)
=
e^{-\int\frac{M_y-N_x}{M}\,dy}
}.
\]

% ============================================================
\subsection{Example: Finding an Integrating Factor}
% ============================================================

\begin{examplebox}{Non-Exact Equation}

Consider

\[
(3xy+y^2)
+
(x^2+xy)\frac{dy}{dx}
=0.
\]

Identify

\[
M(x,y)=3xy+y^2,
\qquad
N(x,y)=x^2+xy.
\]

First,

\[
M_y=3x+2y
\]

and

\[
N_x=2x+y.
\]

Since

\[
M_y\neq N_x,
\]

the equation is \textbf{not exact}.

\end{examplebox}

We look for an integrating factor depending only on $x$.

Thus,

\[
\frac{d\mu}{\mu}
=
\frac{M_y-N_x}{N}\,dx.
\]

Substitute:

\[
\frac{d\mu}{\mu}
=
\frac{(3x+2y)-(2x+y)}
{x^2+xy}\,dx.
\]

Simplifying,

\[
\frac{d\mu}{\mu}
=
\frac{x+y}{x^2+xy}\,dx.
\]

Factor the denominator:

\[
x^2+xy=x(x+y).
\]

Therefore,

\[
\frac{d\mu}{\mu}
=
\frac{x+y}{x(x+y)}\,dx
=
\frac{1}{x}\,dx.
\]

Integrating,

\[
\int\frac{d\mu}{\mu}
=
\int\frac{1}{x}\,dx.
\]

Hence,

\[
\ln|\mu|
=
\ln|x|+C.
\]

Up to a nonzero constant multiple,

\[
\boxed{\mu(x)=x}.
\]

% ============================================================
\subsubsection{Making the Equation Exact}
% ============================================================

Multiply the original equation by $\mu=x$:

\[
x(3xy+y^2)
+
x(x^2+xy)\frac{dy}{dx}
=0.
\]

Thus,

\[
(3x^2y+xy^2)
+
(x^3+x^2y)\frac{dy}{dx}
=0.
\]

Define

\[
\widetilde{M}(x,y)
=
3x^2y+xy^2
\]

and

\[
\widetilde{N}(x,y)
=
x^3+x^2y.
\]

Now,

\[
\widetilde{M}_y
=
3x^2+2xy
\]

and

\[
\widetilde{N}_x
=
3x^2+2xy.
\]

Therefore,

\[
\widetilde{M}_y
=
\widetilde{N}_x,
\]

so the new equation is exact.

% ============================================================
\subsubsection{Finding the Solution}
% ============================================================

Since

\[
\psi_x
=
3x^2y+xy^2,
\]

integrate with respect to $x$:

\[
\psi(x,y)
=
\int(3x^2y+xy^2)\,dx+h(y).
\]

Thus,

\[
\psi(x,y)
=
x^3y+\frac{x^2y^2}{2}+h(y).
\]

Differentiate with respect to $y$:

\[
\psi_y
=
x^3+x^2y+h'(y).
\]

Since

\[
\psi_y
=
\widetilde{N}
=
x^3+x^2y,
\]

we obtain

\[
h'(y)=0.
\]

Therefore,

\[
h(y)=C.
\]

Hence,

\[
\boxed{
x^3y+\frac{x^2y^2}{2}=C
}.
\]

% ============================================================
\subsection{Integrating Factor Summary}
% ============================================================

\begin{formulabox}

For

\[
M(x,y)+N(x,y)\frac{dy}{dx}=0,
\]

if the equation is not exact,

\[
M_y\neq N_x,
\]

check

\[
\boxed{
\frac{M_y-N_x}{N}
}
\]

for dependence on $x$ alone.

If so,

\[
\boxed{
\mu(x)
=
e^{\int\frac{M_y-N_x}{N}\,dx}
}.
\]

Otherwise, check

\[
\boxed{
-\frac{M_y-N_x}{M}
}
\]

for dependence on $y$ alone.

If so,

\[
\boxed{
\mu(y)
=
e^{-\int\frac{M_y-N_x}{M}\,dy}
}.
\]

\end{formulabox}

% ============================================================
\section{Summary}
% ============================================================

\begin{keyideabox}

For a differential equation

\[
M(x,y)+N(x,y)\frac{dy}{dx}=0,
\]

the first question is whether it is exact.

\[
\boxed{M_y=N_x}
\]

If yes, find the potential function $\psi$ and solve

\[
\boxed{\psi(x,y)=C}.
\]

If not, look for an integrating factor $\mu$.

\end{keyideabox}

% ============================================================
\subsection{Exact Equations: Standard Procedure}
% ============================================================

\textbf{01. IDENTIFY}

Identify $M(x,y)$ and $N(x,y)$:

\[
M(x,y)+N(x,y)\frac{dy}{dx}=0.
\]

\textbf{02. CHECK EXACTNESS}

Compute

\[
M_y
\qquad\text{and}\qquad
N_x.
\]

Check whether

\[
\boxed{M_y=N_x}.
\]

\textbf{03. FIND THE POTENTIAL FUNCTION}

If the equation is exact, use

\[
\psi_x=M.
\]

Integrate with respect to $x$:

\[
\psi(x,y)
=
\int M(x,y)\,dx+h(y).
\]

\textbf{04. DETERMINE $h(y)$}

Differentiate with respect to $y$ and use

\[
\psi_y=N.
\]

Solve for $h'(y)$ and integrate.

\textbf{05. WRITE THE SOLUTION}

Set the potential function equal to a constant:

\[
\boxed{\psi(x,y)=C}.
\]

% ============================================================
\subsection{Integrating Factors: Standard Procedure}
% ============================================================

If

\[
M_y\neq N_x,
\]

the equation is not exact.

\textbf{01. CHECK FOR $\mu(x)$}

Compute

\[
\frac{M_y-N_x}{N}.
\]

If this depends only on $x$, use

\[
\boxed{
\mu(x)
=
e^{\int\frac{M_y-N_x}{N}\,dx}
}.
\]

\textbf{02. CHECK FOR $\mu(y)$}

If the previous expression does not depend only on $x$, check

\[
-\frac{M_y-N_x}{M}.
\]

If this depends only on $y$, use

\[
\boxed{
\mu(y)
=
e^{-\int\frac{M_y-N_x}{M}\,dy}
}.
\]

\textbf{03. MULTIPLY BY THE INTEGRATING FACTOR}

Multiply the entire differential equation by $\mu$.

Re-identify

\[
\widetilde{M}
\qquad\text{and}\qquad
\widetilde{N}.
\]

\textbf{04. VERIFY EXACTNESS}

Check

\[
\widetilde{M}_y=\widetilde{N}_x.
\]

\textbf{05. FIND THE POTENTIAL FUNCTION}

Find $\psi(x,y)$ using

\[
\psi_x=\widetilde{M}.
\]

Then use

\[
\psi_y=\widetilde{N}
\]

to determine the remaining function.

\textbf{06. WRITE THE SOLUTION}

Finally,

\[
\boxed{\psi(x,y)=C}.
\]

% ============================================================
\subsection{Big Picture}
% ============================================================

\begin{keyideabox}

\[
\boxed{
\begin{array}{c}
M+N\dfrac{dy}{dx}=0
\\[0.5em]
\downarrow
\\[0.5em]
\text{Check }M_y=N_x
\\[0.5em]
\downarrow
\\[0.5em]
\begin{cases}
\text{Exact}
&
\longrightarrow
\psi
\longrightarrow
\psi=C
\\[0.7em]
\text{Not exact}
&
\longrightarrow
\mu
\longrightarrow
\text{Exact}
\longrightarrow
\psi
\longrightarrow
\psi=C
\end{cases}
\end{array}
}
\]

\end{keyideabox}

% ============================================================


\end{document}
